% !TEX root = main.tex
\clearpage
\section{Проверка и анализ результатов}

Во всех 112 измерениях получено точное поэлементное совпадение с эталоном. В журнале размещения подтверждено, что MPI-процессы выполнялись на разных контейнерных узлах. При $N=257$ ранг 0 обрабатывал 129 строк, ранг 1~--- 128, что проверяет сбор результата при неравных размерах блоков.

В таблице~\ref{tab:time} приведены медианы семи измерений. Для $N=257$, $512$ и $2048$ наименьшее время получено с OpenMP. При $N=1024$ наименьшая медиана получена в гибридном режиме.

\begin{table}[H]
\centering
\small
\caption{Медианное время выполнения, мс}\label{tab:time}
\begin{tabular}{@{}rrrrr@{}}
\toprule
$N$ & Последовательно & OpenMP & MPI & MPI + OpenMP \\
\midrule
\TimeRows
\bottomrule
\end{tabular}
\end{table}

На рисунке~\ref{fig:speedup} показано ускорение относительно последовательного варианта. Добавление второго потока на ранг увеличило ускорение MPI для всех четырёх размеров при том же объёме передаваемых матриц.

\begin{figure}[H]
\centering
\begin{tikzpicture}
\begin{axis}[
    width=0.94\textwidth,
    height=6cm,
    xlabel={Размер матрицы $N$},
    ylabel={Ускорение $S$},
    xtick=data,
    xticklabel style={/pgf/number format/1000 sep={}},
    ymin=0,
    ymajorgrids,
    grid style={gray!25},
    legend style={at={(0.5,0.06)},anchor=south,legend columns=3,font=\small,draw=none},
    table/col sep=comma,
]
\addplot[blue!70!black,thick,mark=*] table[x=n,y=openmp] {build/speedup.csv};
\addlegendentry{OpenMP}
\addplot[orange!80!black,thick,dashed,mark=square*] table[x=n,y=mpi] {build/speedup.csv};
\addlegendentry{MPI}
\addplot[teal!75!black,thick,dashdotted,mark=triangle*] table[x=n,y=hybrid] {build/speedup.csv};
\addlegendentry{MPI + OpenMP}
\end{axis}
\end{tikzpicture}
\caption{Ускорение относительно последовательного варианта}\label{fig:speedup}
\end{figure}

При $N=1024$ медиана гибридного варианта составила 50,598~мс, а OpenMP~--- 52,964~мс. Диапазоны 50,200--51,712~мс и 49,174--54,998~мс пересекаются. Устойчивое преимущество гибридного варианта на этом размере по одной серии установить нельзя.

При $N=2048$ в режиме MPI получено 867,074~мс, в гибридном~--- 462,048~мс. Сокращение времени составило $46{,}7\,\%$. Для OpenMP с четырьмя потоками получено 444,171~мс, то есть гибридный запуск с тем же общим числом потоков занял примерно на $4{,}0\,\%$ больше времени. По формуле~\eqref{eq:metrics} эффективность составила $94{,}1\,\%$ для OpenMP, $96{,}4\,\%$ для MPI и $90{,}5\,\%$ для гибридного режима. Более высокая эффективность MPI получена при меньшем числе исполнителей.

Время обменов входит в измерения MPI. Разность между временами OpenMP и гибридного варианта включает также влияние планирования потоков и доступа к памяти.
